\documentclass[portrait,a0,final,20pt]{a0poster}
%%%Load packages
\usepackage{multicol} 			%2-column layout
% \setlength\columnsep{8pt} % This is the default columnsep for all pages
\usepackage[left=2cm,right=2cm,bottom=2cm,top=0cm]{geometry}			%Reset margins
\usepackage{helvet}				%Load Helvetica font & CM math
\usepackage{color}				%Needed for colour boxes & coloured text
% \usepackage{graphics}
\usepackage{graphicx}
\usepackage{subcaption}
\usepackage[font=Large,labelfont=bf]{caption}
\usepackage{float}
\usepackage{tikz}
\usetikzlibrary{arrows,positioning, shapes.symbols,shapes.callouts,patterns,shapes,chains,calc,backgrounds,fadings}

\usepackage[absolute,overlay]{textpos} % for logo next ti title

\usepackage{titlesec}
\usepackage{tabularx, booktabs} % make width of table columns evenly distributed (see http://tex.stackexchange.com/questions/60601/evenly-distributing-column-widths)

\usepackage{enumitem}
\usepackage{animate}[2017/05/18]

\usepackage[hidelinks, colorlinks = true,
            linkcolor = blue,
            urlcolor  = blue,
            citecolor = blue,
            anchorcolor = blue]{hyperref}
\usepackage{xcolor}
\usepackage{newtxtext,newtxmath}

\usepackage[noend]{algorithm2e}
\newcommand\mycommfont[1]{\footnotesize\ttfamily\textcolor{blue}{#1}}
\SetCommentSty{mycommfont}

\usepackage{amssymb} 
\usepackage{pifont}% for red cross and green tick

% \usepackage{titlepic}
% \usepackage{columns}

\titlespacing\section{0pt}{10pt plus 4pt minus 2pt}{5pt plus 2pt minus 2pt}
\titlespacing\subsection{0pt}{12pt plus 4pt minus 2pt}{0pt plus 2pt minus 2pt}

\titleformat*{\section}{\Huge\bfseries}
\titleformat*{\subsection}{\Large\bfseries}
\titleformat*{\subsubsection}{\large\bfseries}
\titleformat*{\paragraph}{\large\bfseries}
\titleformat*{\subparagraph}{\large\bfseries}

%%%Define colours and lengths
\definecolor{headingcol}{rgb}{0,0,0}			%Colour of main title
\definecolor{boxcol}{rgb}{0.7,0.2,0.2}		%Edge-colour of box and top banner
\fboxsep=1cm							%Padding between box and text
\setlength{\columnsep}{2cm}				%Set spacing between columns
\renewcommand{\familydefault}{\sfdefault}	%Set main text to sans-serif

%%%Format title
\makeatletter							%Needed to include code in main file
\renewcommand\@maketitle{%
\null									%Sets position marker
% \vspace{5em}
{
\color{headingcol}\sffamily\VeryHuge \textbf		%Set title font and colour
\@title \par}%
% \vskip 0.6em%
{
\color{black}\sffamily\LARGE				%Set author font and colour
\lineskip .5em%
\begin{tabular}[t]{l}%
\@author
\end{tabular}\par}%
\vskip 1cm
\par
}
\makeatother



\title{Disease Knowledge Transfer across Neurodegenerative Diseases}

\newcommand{\inst}[1]{$^{#1}$}

% \author{\Large{R\u{a}zvan V. Marinescu\inst{1}, Neil P. Oxtoby\inst{1}, Alexandra L. Young\inst{1}, Arman Eshaghi\inst{2}, Peter A. Wijeratne\inst{1}, Daniel C. Alexander\inst{1}}\\
% \begin{tabular}{l p{1cm} l p{1cm} l}
% \large{$^1$Centre for Medical Image Computing, University College London}  & & \large{$^2$Queen Square MS Centre, UCL Institute of Neurology} \\
% \end{tabular}
% }

\author{\LARGE{R\u{a}zvan Marinescu\inst{1,2}, Marco Lorenzi\inst{3}, Stefano Blumberg\inst{1}, Alexandra Young\inst{1}, Pere Morell\inst{1}},\\ \LARGE{Neil Oxtoby\inst{1}, Arman Eshaghi\inst{1,4}, Keir Young\inst{5}, Sebastian Crutch\inst{5}, Polina Golland\inst{2}, Daniel Alexander\inst{1}}\\
\begin{tabular}{l p{1cm} l p{1cm} l}
\large{$^1$Centre for Medical Image Computing, UCL, UK}  & & \large{$^3$University of C\^{o}te d'Azur, Inria Sophia Antipolis, France} & & \large{$^5$Dementia Research Centre, UCL, UK}\\
\large{$^2$Computer Science and Artificial Intelligence Laboratory, MIT, USA} & & \large{$^4$Queen Square MS Centre, UCL Institute of Neurology, UK} & & 
 \\
\end{tabular}
} 

% \institute{Centre for Medical Image Computing, University College London, UK
% \and 
% Computer Science and Artificial Intelligence Laboratory, MIT, USA
% \email{razvan@csail.mit.edu}
% \and
% Queen Square MS Centre, UCL Institute of Neurology, UK
% \and 
% Dementia Research Centre, University College London, UK
% \and
% University of C\^{o}te d'Azur, Inria Sophia Antipolis, France
% % Laboratoire d'Analyse Num\'{e}rique, B\^{a}timent 425,\\
% % F-91405 Orsay Cedex, France}
% }

\newcommand{\fnt}[1]{\LARGE{#1}}

\DeclareMathOperator*{\argmin}{arg\,min}
\DeclareMathOperator*{\argmax}{arg\,max}


\begin{document}
% \hspace{-3cm}								%Align with edge of page, not margin
% \vspace{2cm}
% \includegraphics[scale=1,trim=0 0 0 200, clip]{Black_Landscape.pdf}

\begin{minipage}{50cm}					%Minipage for title contents
% \vspace{-18cm}
\maketitle
\end{minipage}
\begin{textblock}{0}(12.2,0.2)
\begin{tabular}{l}
\includegraphics[height=5.5cm]{mit_logo}\includegraphics[height=6cm]{qr_dkt.png}\\
\includegraphics[height=5cm]{ucl_logo}\\
\end{tabular}


\end{textblock}

\fnt{


\newcommand{\lp}{\lambda_{d_i}^{\psi(k)}}
\newcommand{\lpuu}{\lambda_{d_i}^{\psi(k),(u)}}
\newcommand{\lpum}{\lambda_{d_i}^{\psi(k),(u-1)}}

\newcommand{\cmark}{\ding{51}}%
\newcommand{\xmark}{\ding{55}}%

\newcommand{\yes}{{\LARGE \textcolor{green!50!black}{\cmark} \par}}
\newcommand{\no}{{\LARGE \textcolor{red}{\xmark} \par}}

\hrule

\newcommand{\heading}[1]{\noindent\textbf{\Huge{#1}}}

\begin{multicols}{2}							
\raggedcolumns							%Don't stretch contents vertically

\pagenumbering{gobble}

%%%Column1
\vspace{-3em}

\heading{Aim} Infer progression of non-MRI biomarkers in rare neurodegenerative diseases (NDs) by leveraging larger datasets of common NDs.


\heading{Why} Posterior Cortical Atrophy (PCA): progression of non-MRI biomarkers not known $\rightarrow$ Identify outcome measures and suitable subjects for PCA clinical trials

\columnbreak

\heading{Challenges}



% \begin{itemize}
%  \item Datasets on rare neurodegenerative diseases (e.g., Posterior Cortical Atrophy) are unimodal (MRI only), cross-sectional and small.
%  \item The continuous progression of non-MRI markers in rare neurodegenerative diseases is not well understood
% \end{itemize}


\begin{figure}[H]
\fnt{
\begin{subfigure}{0.48\columnwidth}
  Typical Neurodeg. Diseases
  
  \begin{itemize}
   \item Large datasets \yes
   \item Multimodal imaging \yes
   \item Longitudinal \yes
%    \item Well-understood \yes
   
  \end{itemize}
\end{subfigure}
\begin{subfigure}{0.48\columnwidth}
 
   Rare Neurodeg. Diseases
  
  \begin{itemize}
   \item Small datasets \no
   \item MRI only \no
   \item Cross-sectional \no
%    \item Poorly understood \no
  \end{itemize}
 
\end{subfigure}
 }
\end{figure}

\end{multicols}
\hrule

\begin{multicols}{2}							
\raggedcolumns	

\heading{Method}
 
\vspace{0.6em}
\begin{figure}[H]
 \centering
   \begin{subfigure}{0.49\columnwidth}
   \centering
   \fnt{1. Each disease characterised by region-specific dysfunction profile}\\
   $ \gamma_{ij}^l = f(\beta_{i} + m_{ij}; \lambda_{d_i}^l) $\\
%    \vspace{0.9em}
%    \includegraphics[width=\textwidth]{diseaseEvol}
   
  \end{subfigure}
  \begin{subfigure}{0.49\columnwidth}
   \centering
%    \vspace{-1.2em}
   \fnt{2. Dysfunction score modelled using region-specific biomarkers}\\
   $ y_{ijk} = g( \gamma_{ij}^{\psi(k)} ; \theta_k) + N(0,\epsilon_k) $
%    \vspace{0.6em}
%    \includegraphics[width=\textwidth]{dysfuncEvol}
  \end{subfigure}

 \includegraphics[width=\columnwidth]{disease_knowledge_transfer_poster.pdf}
  \fnt{3. Extend to multiple subjects, biomarkers and diseases}\\
 $ p(\boldsymbol{y}|\theta, \lambda, \beta , \epsilon) = \prod_{(i,j,k) \in \Omega} p(y_{ijk}|\theta_k, \lp, \beta_{i}) $
 \end{figure}



\heading{Inference:} Loopy belief propagation

% TODO bold out the diferences

\newcommand{\uu}{^{(u)}}
\newcommand{\um}{^{(u-1)}}

\newcommand{\algoFnt}[1]{\Large{#1}}


\begin{figure}[H]
\begin{algorithm}[H]
\algoFnt
 Initialise $\boldsymbol{\theta}^{(0)}$, $\boldsymbol{\lambda}^{(0)}$, $\boldsymbol{\beta}^{(0)}$\\
  \While{$\boldsymbol{\theta}$, $\boldsymbol{\lambda}$, $\boldsymbol{\beta}$ not converged}{
   \tcp*[l]{\algoFnt{Estimate biomarker trajectories (disease agnostic)}}
%     \For{$k=1$ to $K$}{
      ${\theta_k\uu = \argmin_{\theta_k} \textcolor{red}{\sum_{(i,j) \in \Omega_k}} \left[y_{ijk} - g\left(f(\beta_i\um + m_{ij}; \lpum) ; \theta_k\right) \right]^2 }$\\
%       ${\epsilon_k\uu = \frac{1}{|\Omega_k|} \sum_{(i,j) \in \Omega_k}    \left[y_{ijk} - g\left(f(\beta_i\um + m_{ij}; \lpum) ; \theta_k\uu \right) \right]^2 }$\\
%     }
     \tcp*[l]{\algoFnt{Estimate dysfunction trajectories (disease specific)}} 
%     \For{$d=1 \in \mathbb{D}$}{
%       \For{$l=1 \in \Lambda$}{
        ${\lambda_{d}^{l, (u)} = \argmin_{\lambda_{d}^{l}} \textcolor{red}{\sum_{(i,j,k) \in \Omega_{d,l}}} \left[y_{ijk} - g\left(f(\beta_i\um + m_{ij}; \lambda_{d}^{l}) ; \theta_k\uu 
        \right) \right]^2 }$\\
%       }
%     }
    \tcp*[l]{\algoFnt{Estimate subject-specific time shifts}} 
%     \For{$i=1 \in [1, \dots, S]$}{
      ${\beta_i\uu = \argmin_{\beta_i} \textcolor{red}{\sum_{(j,k) \in \Omega_i}} \left[y_{ijk} - g\left(f(\beta_i + m_{ij}; \lpuu) ; \theta_k\uu
      \right) \right]^2 }$\\
%     }
}
% \normalfont
\end{algorithm}
% \caption[The algorithm for estimating the DKT parameters]{The algorithm used to estimate the DKT parameters, based on loopy belief-propagation.}
\label{fig:dktAlgo}
\end{figure}

\heading{Demographics}

\begin{itemize}
 \item Dementia Research Center cohort: 76 PCA, 67 tAD, 87 age-matched controls for training, 10 PCA with DTI for validation
 \item TADPOLE Challenge dataset split into three subgroups with different progressions: 21 hippocampal, 35 cortical, 27 subcortical.
 \item Synthetic dataset mimicking the DRC cohort: 50 subjects with "synthetic PCA", 100 subjects with "synthetic AD".
\end{itemize}


% The cohort from the Dementia Research Centre UK used for validation consisted of 10 subjects diagnosed with Posterior Cortical Atrophy, with a mean age of 59.4, 40\% females, as well as 10 age-matched controls with a mean age of 59.3, 50\% females.
% 
% 
% For the validation on TADPOLE subgroups, we used applied the SuStaIn model on TADPOLE to split the population into three subgroups with different progression: hippocampal, cortical and subcortical subypes with prominent atrophy in the hippocampus, cortical and subcortical areas respectively.
% The resulting subgroups had the following demographics:
% 
% \begin{table}
%  \begin{tabular}{c | c c c c}
% \textbf{Cohort} & \textbf{Nr. subjects} & \textbf{Nr. visits} & \textbf{Age (baseline)} & \textbf{Gender (\%F)}\\
% \hline
% Controls (Hippocampal) & 31 & 2.3 $\pm$ 1.8 & 74.4 $\pm$ 6.9 & 38\%\\
% AD (Hippocampal) & 21 & 1.5 $\pm$ 0.8 & 74.5 $\pm$ 5.5 & 42\%\\
% \hline
% Controls (cortical) & 21 & 2.3 $\pm$ 1.3 & 70.9 $\pm$ 5.4 & 42\%\\
% AD (cortical) & 35 & 1.7 $\pm$ 0.9 & 72.8 $\pm$ 7.4 & 28\%\\
% \hline
% Controls (subcortical) & 28 & 3.0 $\pm$ 1.5 & 73.7 $\pm$ 6.5 & 42\%\\
% AD (subcortical) & 27 & 1.6 $\pm$ 0.9 & 73.7 $\pm$ 7.5 & 33\%\\
%  \end{tabular}
%  \caption{Demographics of the subjects in the three TADPOLE subgroups.}
%  \label{tab:demogTadSubtypes}
% \end{table}

\columnbreak

\heading{Results}

\vspace{0.6em}
\begin{figure}[H]
\fnt{In synthetic experiment, the estimated parameters are close to the true parameters}\\
\includegraphics[width=\columnwidth]{../figures/compTrueParams105_synth1_JMD.pdf}
  \label{fig:dktSynthTrajCompTrue}
\end{figure}

\begin{figure}[H]
% \centering
\fnt{Inferred multimodal trajectories for Posterior Cortical Atrophy are plausible, suggesting late-stage parietal and occipital damage. }
 \includegraphics[width=\columnwidth, trim=0 20 0 35, clip]{../figures/trajDisSpaceOverlap_PCA_tad-drcTinyPen5_JMD.png}
%  \caption{Estimated trajectories for the PCA cohort. The only data that were available were the MRI volumetric data. The dynamics of the other biomarkers has been inferred by the model using data from typical AD, and taking into account the different spatial distribution of pathology in PCA vs tAD. }
%  \label{fig:PCAtrajByModality}
\end{figure}


\noindent \fnt{Our model has favourable performance compared to other models, on two different datasets}
\vspace{0.3em}
\begin{figure}[H]
\fontsize{27}{30}\selectfont
\begin{tabular}{c | c c c c c c}
\textbf{Model} & \textbf{Cingulate} & \textbf{Frontal} & \textbf{Hippocam.} & \textbf{Occipital} & \textbf{Parietal} & \textbf{Temporal}\\
& \multicolumn{6}{c}{\textbf{TADPOLE: Hippocampal subgroup to Cortical subgroup}}\\
DKT (ours) &      0.56 $\pm$ 0.23 &    \textbf{0.35 $\pm$ 0.17} &        \textbf{0.58 $\pm$ 0.14} &     -0.10 $\pm$ 0.29 &     \textbf{0.71 $\pm$ 0.11} &     \textbf{0.34 $\pm$ 0.26} \\
Latent stage &      0.44 $\pm$ 0.25 &    0.34 $\pm$ 0.21 &       0.34 $\pm$ 0.24* &     \textbf{-0.07 $\pm$ 0.22} &     0.64 $\pm$ 0.16 &    0.08 $\pm$ 0.24* \\
Multivariate &      \textbf{0.60 $\pm$ 0.18} &   0.11 $\pm$ 0.22* &       0.12 $\pm$ 0.29* &     -0.22 $\pm$ 0.22 &   -0.44 $\pm$ 0.14* &   -0.32 $\pm$ 0.29* \\
Spline &    -0.24 $\pm$ 0.25* &  -0.06 $\pm$ 0.27* &        0.58 $\pm$ 0.17 &     -0.16 $\pm$ 0.27 &    0.23 $\pm$ 0.25* &    0.10 $\pm$ 0.25* \\
Linear &    -0.24 $\pm$ 0.25* &   0.20 $\pm$ 0.25* &        0.58 $\pm$ 0.17 &     -0.16 $\pm$ 0.27 &    0.23 $\pm$ 0.25* &    0.13 $\pm$ 0.23* \\
& \multicolumn{6}{c}{\textbf{typical Alzheimer's to Posterior Cortical Atrophy}}\\
DKT (ours) &    0.77 $\pm$ 0.11 &    0.39 $\pm$ 0.26 &      0.75 $\pm$ 0.09 &    0.60 $\pm$ 0.14 &    \textbf{0.55 $\pm$ 0.24} &    \textbf{0.35 $\pm$ 0.22} \\
Latent stage &    \textbf{0.80 $\pm$ 0.09} &    \textbf{0.53 $\pm$ 0.17} &      \textbf{0.80 $\pm$ 0.12} &    0.56 $\pm$ 0.18 &    0.50 $\pm$ 0.21 &    0.32 $\pm$ 0.24 \\
Multivariate &   0.73 $\pm$ 0.09 &   0.45 $\pm$ 0.22  &    0.71 $\pm$ 0.08 & -0.28 $\pm$ 0.21* &  0.53 $\pm$ 0.22  &  0.25 $\pm$ 0.23* \\
Spline &   0.52 $\pm$ 0.20* &  -0.03 $\pm$ 0.35* &     0.66 $\pm$ 0.11* &   0.09 $\pm$ 0.25* &    0.53 $\pm$ 0.20 &   0.30 $\pm$ 0.21* \\
Linear &   0.52 $\pm$ 0.20* &    0.34 $\pm$ 0.27 &     0.66 $\pm$ 0.11* &    \textbf{0.64 $\pm$ 0.17} &    0.54 $\pm$ 0.22 &   0.30 $\pm$ 0.21* \\
\end{tabular}
\vspace{0.5em}
% \caption[Performance evaluation of DKT and other models]{Performance evaluation of DKT and four other statistical models of decreasing complexity. }
% \end{footnotesize}
\label{sec:dktPerfMetrics}
\end{figure}


\heading{Conclusion}

\begin{itemize}
 \item Developed a novel methodology and model for transfer learning across different diseases
 \item Inferred multimodal trajectories for Posterior Cortical Atrophy
%  \item 
\end{itemize}





\end{multicols}
\hrule


}


\begin{multicols}{2}

\Large{

\raggedcolumns	


\vspace{-1em}
\subsection*{Weblinks}
\begin{itemize}
\item Source code: \url{https://github.com/mrazvan22/dkt}
\item Website: \url{https://people.csail.mit.edu/razvan/}
\end{itemize}
}



\columnbreak

\subsection*{Funders}

% \begin{figure}[H]
% \begin{subfigure}{0.48\columnwidth}
\vspace{0em}
% \hspace{2em}
\newcommand{\heightLogos}{3.5cm}
\includegraphics[height=\heightLogos]{epsrc_logo}
\hspace{0.5em}
\includegraphics[height=\heightLogos]{nac_logo} 
\hspace{0.5em}
\includegraphics[height=\heightLogos]{cdt_logo} 
\hspace{0.5em}
\includegraphics[height=\heightLogos]{nih_logo} 
\hspace{0.5em}
\includegraphics[height=\heightLogos]{aruk_logo} 
\hspace{0.5em}
% \end{subfigure}
% \begin{subfigure}{0.48\columnwidth}
%  \begin{itemize}
%  \item  EPSRC grant EP/L016478/1
%  \item  NIH grant NIBIB NAC P41EB015902
% \end{itemize}
% \end{subfigure}

% \end{figure}



\end{multicols}


\end{document}
